% Source: https://artofproblemsolving.com/community/c3h2186663_interesting_problem
% Stable identity: aops-topic-2186663-q01
Let each payment be $x$ dollars and let $r=1.06$.

\textbf{(a)} After the first year I owe $10000r-x$; after the second, $(10000r-x)r-x$; after the third,
\[10000r^3-xr^2-xr-x=0 .\]
So
\[x=\dfrac{10000r^3}{r^2+r+1}=\dfrac{10000\cdot1.191016}{1.1236+1.06+1}=\dfrac{11910.16}{3.1836}\approx3741.098,\]
and each payment is $\boxed{\$3741.10}$.

\textbf{(b)} The present value of a payment is the amount of money which, invested today at the same $6\%$ interest, would grow to that payment by the time it is made. The first payment is made at the end of the first year, \emph{after} the first year's interest has been added; an amount $y$ invested today grows to $1.06y$ by then, so the present value of the first payment is $y=\dfrac{x}{1.06}$. (It is not $x$: by the time I pay, the debt has already been charged one year of interest.) Similarly the second and third payments have present values $\dfrac{x}{1.06^2}$ and $\dfrac{x}{1.06^3}$.

Dividing the equation in (a) by $r^3$ gives
\[10000=\dfrac{x}{r}+\dfrac{x}{r^2}+\dfrac{x}{r^3},\]
which says exactly that the present values of the three payments add up to the amount borrowed. This gives the same payment: $x=\dfrac{10000}{\dfrac{1}{1.06}+\dfrac{1}{1.06^2}+\dfrac{1}{1.06^3}}\approx3741.10$.
